%Text automatically generated by txt2latex
\section{NAME}
txt2latex - convert flat ASCII text to \LaTeX.

\section{SYNOPSIS}
txt2latex [OPTION...] FILE 

\section{DESCRIPTION}
txt2latex converts the input text into \LaTeX. 

txt2latex is also able to recognize and format sections, paragraphs,
lists (itemize, enumerate, description) and verbatim blocks.

If input file FILE is omitted, standard input is used. 
Result is displayed on standard output. 

Here is how text patterns are recognized and processed:
\begin{description}
\item[Sections] These headers are defined by a line in upper case, starting
                column 1. 
                Optionally, the Section name can be preceded by a blank line. 
                This is useful for a better visualization of the source 
                text to be used to generate the \LaTeX source code.
\item[Paragraphs] They must be separated by a blank line, and left aligned.
                  Alternatively two blank spaces can be used to produce the
                  same result. This option will provide a better visualization
                  of the source text to be used to generate the \LaTeX source code.
\item[Description list] 
                        The item definition is separated from the item description
                        by at least 2 blank spaces, even before a new line, if
                        definition is too long.
\item[Bullet list] 
                   Bullet list items are defined by the first word being "-",
                   "*" or "o".
\item[Enumerated list] 
                       The first word must be a number followed by a dot or a rounded bracket.
\item[Verbatim blocks] 
                       This paragraph type is used to display unmodified text,
                       for example source code. It must be separated by a blank
                       line and be indented by a TAB. It is primarily used to format
                       unmodified source code. It will be printed using verbatim environment.
\item[Mathematics] 
                   Inline mathematics must be enclosed with a simple \$ sign
                   and equations must be enclosed with double \$ signs.
\end{description}
              
For example, $E=mc^2$ is rendered as an inline math whereas 
$$ E=mc^2 $$
is rendered in a simple equation environment.
   

\section{OPTIONS}
\begin{description}
\item[-v, --version] Display version.
\item[-h, --help] Display help.
\item[-d date] Set date. Defaults to current date.
\item[-t mytitle] Set the title. If the title is set, txt2latex will 
                  automatically add the preambule and markups for the document
\item[-a author] Set the author.
\item[-s shift] Shift heading level by 0 (section), 1 (subsection), or 2 (subsection).
                Defaults to 0.
\item[-I txt] Italicize txt in output. Can be specified more than once.
\item[-B txt] Emphasize (bold) txt in output. Can be specified more than once.
\item[-P package] Add packages or LateX commands in the preambule.
\item[-X] Compile output with pdflatex.
\end{description}

\section{NOTES}
The formatting rules are heavily inspired from txt2man(1). 
The special treatment for sections NAME and SYNOPSIS performed 
automatically by txt2man(1) is not done by txt2latex.
Nonetheless, the objective of txt2latex is not to parse on man-formatted
plain text but to convert any plain text to \LaTeX. 

Using the options -B and -I, you can easily format your text as it 
would be done by txt2man. 

\section{EXAMPLES}
Simple conversion

\begin{verbatim}
    $ txt2latex FILE > FILE.tex
  
\end{verbatim}
Conversion with document title

\begin{verbatim}
    $ txt2latex -t "title" -a "author" -I "word" FILE > FILE.tex

\end{verbatim}
Conversion with custom packages and direct rendering with pdflatex

\begin{verbatim}
    $ txt2latex -t "title" -a "author" -I "word" -P "\\usepackage{ccfonts}" -X FILE

\end{verbatim}
\section{SEE ALSO}
txt2man(1)
